% Source-derived chapter 09: Special cycles of corank one
% Original source: higher-siegel-weil-unitary-nonsingular-terms
\chapter{Special cycles of corank one}
\label{higher-siegel-weil-unitary-nonsingular-terms-chapter-09}
\source{higher-siegel-weil-unitary-nonsingular-terms}

\begin{remark}[Source section]
\label{higher-siegel-weil-unitary-nonsingular-terms-chapter-09-source}
\source{higher-siegel-weil-unitary-nonsingular-terms}
\source{higher-siegel-weil-unitary-nonsingular-terms}
This chapter reproduces the corresponding section of the retrieved
LaTeX source, including its definitions, statements, examples, and
proofs. It has not yet been translated into Lean.
\notready
\end{remark}

% The source section title was: Special cycles of corank one
\label{ssec: geometric properties}
\source{higher-siegel-weil-unitary-nonsingular-terms}

In this section we prove geometric properties of the special cycles $\cZ^{r}_{\cE}(a)$ when $m= \rank \cE = 1$ (where the number field analogues are called ``Kudla-Rapoport divisors''). In particular, we show that for $a\ne0$, $\cZ^{r}_{\cE}(a)$  are local complete intersections of dimension $r$ less than $\Sht^{r}_{U(n)}$.  When $a=0$,  we show that  $\cZ^{r}_{\cE}(a)^{*}$ has dimension at most $\dim\Sht^{r}_{U(n)}-r$. These geometric properties are proved by studying stratifications introduced and analyzed in \S \ref{ss:strat KR div} and \S \ref{ssec: m=1 a=0 case}.

%In this section we prove geometric properties of special cycles to justify the definition of the cycle classes in \S\ref{s:int prob}. In particular, we establish the following properties.
%\begin{itemize}
%\item When $m= \rank \cE = 1$ (where the number field analogues are called ``Kudla-Rapoport divisors''), the $\cZ^{r}_{\cE}(a)$ are local complete intersections of dimension $r$ less than $\Sht^{r}_{U(n)}$.
%\item When $ m = \rank \cE = 1$, the cycle class $[\cZ^{r}_{\cE}(a)]$ can be identified with the class defined using the Hitchin moduli stack. 
%\end{itemize} 


%
%\subsection{Stratification of the special cycle when $m=1$}\label{ss:strat KR div}
%Consider the a fixed line bundle $\cL$ on $X'_{k}$ and $a\in \cA_{\cL}(k)$, i.e.,  $a: \cL\to \s^{*}\cL^{\vee}$ is nonzero Hermitian.  For each $(\cL\xr{t_{i}}\cF_{i})_{0\le i\le r}\in \cZ(\ov k)$ with legs $(x'_{i})_{1\le i\le r}\in X'(\ov k)^{r}$, let $D_{i}$ ($0\le i\le r$) be the divisor on $X'_{\ov k}$ such that $t_{i}: \cL(D_{i})\incl \cF_{i}$ is saturated. For each $1\le i\le r$, we have one of the  four cases:
%\begin{enumerate}
%\item[Case (0)] $D_{i}=D_{i-1}$;
%\item[Case (+)] $D_{i}=D_{i-1}+\s(x'_{i})$;
%\item[Case ($-$)] $D_{i}=D_{i-1}-x'_{i}$;
%\item[Case ($\pm$)] $D_{i}=D_{i-1}-x'_{i}+\s(x'_{i})$.
%\end{enumerate}
%Since the composition $\cL\xr{t_{i}}\cF_{i}\xr{h}\s^{*}\cF_{i}^{\vee}\xr{\s^{*}t_{i}^{\vee}}\s^{*}\cL^{\vee}$ is equal to $a$, we see that $D_{i}+\s(D_{i})$ is a subdivisor of the divisor of $a$. Therefore $\nu(D_{i})\le D_{a}$ as divisors on $X(\ov{k})$. 
%
%
%\begin{defn} A {\em saturation type} of length $r$ is a function $\d: \{1,2,\cdots, r\}\to \{0,+,-,\pm\}$ such that $\d^{-1}(-)$ and $\d^{-1}(+)$ have the same cardinality.  Let $\D_{r}$ be the set of saturation types of length $r$. 
%\end{defn}
%
%Let $d_{0}\in \Z_{\ge0}$ and $\d\in\D_{r}$. Let $\frD_{d_{0}}[\d]$ be the moduli space of $(\{x'_{i}\}_{1\le i\le r}, \{D_{i}\}_{0\le i\le r})$ where
%\begin{enumerate}
%\item  $x'_{i}\in X'$ for $1\le i\le r$;
%\item $D_{i}$ are effective divisors on $X'$ for $0\le i\le r$ and $\deg D_{0}=d_{0}$.
%\item  For each $1\le i\le r$, the triple $(D_{i-1}, D_{i},x'_{i})$ falls into one of the case indexed by $\d(i)\in \{0,+,-,\pm\}$ above. 
%\end{enumerate}
%Since $D_{i}$ for $i\ge1$ are uniquely determined by $D_{0}$ and $\{x'_{i}\}_{1\le i\le r}$, we have an embedding $\frD_{d_{0}}[\d]\incl X'_{d_{0}}\times X'^{r}$.
%
%Let $\frM_{d_{0}}$ be the moduli stack classifying $(\cL, D_{0}, (\cF,h), t: \cL\to \cF)$ where $(\cL,(\cF,h),t)\in \cM(1,n)$ (see Definition \ref{}) such that $a(t)=\s^{*}t^{\vee}\c h\c t\ne0$, $D_{0}$ an effective divisors of degree $d_{0}$ on $X'$ such that $t$ extends to a saturated map $t: \cL(D_{0})\to \cF$.
%
%Let $\frH^{r}_{d_{0}}[\d]$ be the moduli stack classifying $(\cL, \{x'_{i}\}_{1\le i\le r}, \{D_{i}\}_{0\le i\le r}, \{\cF_{i},h_{i}\}_{0\le i\le r}, \{t_{i}: \cL\to \cF_{i}\})$ where 
%\begin{enumerate}
%\item the tuple $(\cL, \{x'_{i}\}, \{\cF_{i},h_{i}\}, \{t_{i}\})\in \Hk^{r}_{\cM(1,n)}$ (see Definition \ref{}) with $a=a(t_{i})=\s^{*}t^{\vee}_{i}\c h_{i}\c t_{i}\ne0$.
%\item the tuple $(\{x'_{i}\}_{1\le i\le r}, \{D_{i}\}_{0\le i\le r})\in \frD_{d_{0}}[\d]$. 
%\item Each $t_{i}$ extends to a saturated map $t: \cL(D_{0})\to \cF$.
%\end{enumerate}
%We have maps $p_{0}, p_{r}: \frH^{r}_{d_{0}}[\d]\to \frM_{d_{0}}$ recording $(\cL, D_{0}, \cF_{0}, h_{0}, t_{0})$ and  $(\cL, D_{r}, \cF_{r}, h_{r}, t_{r})$. 
%
%Finally we define the Shtuka version $\frS^{r}_{d_{0}}[\d]$ of $\frH^{r}_{d_{0}}[\d]$ by the Cartesian diagram
%\begin{equation}
%\xymatrix{\frS^{r}_{d_{0}}[\d] \ar[r]\ar[d] & \frH^{r}_{d_{0}}[\d]\ar[d]^{(p_{0},p_{r})}\\
%\frM_{d_{0}}\ar[r]^-{(\Id,\Frob)} & \frM_{d_{0}}\times \frM_{d_{0}}}
%\end{equation}
%Just as $\Sht^{r}_{\cM(m,n)}$ decomposes into the disjoint union of $\cZ^{r}_{\cE}(a)$  for $(\cE,a)\in \cA(m,n)(k)$, we have a decomposition
%\begin{equation}
%\frS^{r}_{d_{0}}[\d]=\coprod_{(\cL,a)\in \cA(1,n)(k)}\cZ^{r}_{\cL}(a)_{d_{0}}[\d]
%\end{equation}
%where $\cZ^{r}_{\cL}(a)_{d_{0}}[\d]$ is the substack of $\frS^{r}_{d_{0}}[\d]$ with fixed $\cL$ and $a=\s^{*}t_{i}^{\vee}\c t_{i}: \cL\to \s^{*}\cL^{\vee}$. From the definitions it is easy to see the following. 
%
%\begin{lemma}
%Let $(\cL,a)\in \cA(1,n)(k)$.
%\begin{enumerate}
%\item Each $\cZ^{r}_{\cL}(a)_{d_{0}}[\d]$ is a locally closed substack of $\cZ^{r}_{\cL}(a)$, and it is empty if $d_{0}>\deg D_{a}$.  
%\item The non-empty substacks in $\{\cZ^{r}_{\cL}(a)_{d_{0}}[\d];0\le d_0\le\deg D_{a}, \d\in\D_{r}\}$ form a finite partition of $\cZ^{r}_{\cL}(a)$. 
%\end{enumerate}
%\end{lemma}


\subsection{Stratification of the special cycle when $m=1$}\label{ss:strat KR div}
\source{higher-siegel-weil-unitary-nonsingular-terms}
In this section we fix a line bundle $\cL$ on $X'_{k}$ and $a\in \cA_{\cL}(k)$, i.e.,  $a: \cL\to \s^{*}\cL^{\vee}$ is nonzero Hermitian. We now define a stratification of $\cZ:=\cZ_{\cL}^{r}(a)_{\ov k}$ and estimate the dimension of each stratum.

When $n=1$, we have $\dim \cZ  =0$. Indeed, Lemma \ref{l:KR cycle legs} implies that there is a $0$-dimensional closed subscheme $D \subset (X')^r$ such that the morphism $\cZ \rightarrow \Sht_{U(1)}^r$ takes values in the preimage of $D$ under the leg map $\Sht_{U(1)}^r \rightarrow (X')^r$. This preimage is $0$-dimensional by Lemma \ref{lem: shtuka geometry} (2). As $Z \rightarrow \Sht_{U(1)}^r$ is finite (by Proposition \ref{prop: Z finite}) with 0-dimensional image, we conclude that $\dim \cZ = 0$. 

Therefore in the rest of this subsection we will assume $n\ge2$. 
 
For each $(\cL\xr{t_{i}}\cF_{i})_{0\le i\le r}\in \cZ(\ov k)$ with legs $(x'_{i})_{1\le i\le r}\in X'(\ov k)^{r}$, let $D_{i}$ ($0\le i\le r$) be the divisor on $X'_{\ov k}$ such that $t_{i}: \cL(D_{i})\incl \cF_{i}$ is saturated. For each $1\le i\le r$, we have one of the  four cases:
\begin{enumerate}
\item[(0)] $D_{i}=D_{i-1}$;
\item[(+)] $D_{i}=D_{i-1}+\s(x'_{i})$;
\item[($-$)] $D_{i}=D_{i-1}-x'_{i}$;
\item[($\pm$)] $D_{i}=D_{i-1}-x'_{i}+\s(x'_{i})$.
\end{enumerate}
Since the composition $\cL\xr{t_{i}}\cF_{i}\xr{h}\s^{*}\cF_{i}^{\vee}\xr{\s^{*}t_{i}^{\vee}}\s^{*}\cL^{\vee}$ is equal to $a$, we see that $D_{i}+\s(D_{i})$ is a subdivisor of the divisor of $a$. Therefore $\nu(D_{i})\le D_{a}$ as divisors on $X(\ov{k})$. 

\subsubsection{Indexing set for strata} Consider the set $\frD$ of sequences of effective divisors $(D_{i})_{0\le i\le r}$ on $X'_{\ov{k}}$ satisfying
\begin{itemize}
\item $\nu(D_{i})\le D_{a}$ for all $0\le i\le r$.
\item For each $1\le i\le r$, the pair $(D_{i-1}, D_{i})$ falls into one of the cases $(0),(+),(-),(\pm)$ above for some $x'_{i}\in X'(\ov k)$.
\item $D_{r}={}^{\t}D_{0}$.
\end{itemize}
It is clear that $\frD$ is a finite set. This will be the index set for our stratification of $\cZ$. 

\subsubsection{Definition of strata} Fix $D_{\bu}=(D_{i})_{0\le i\le r}\in \frD$. Let $I_{0} :=\{1\le i\le r|D_{i}=D_{i-1}\}$. Similarly we define $I_{+}, I_{-}$ and $I_{\pm}$ as the set of those $i$ such that $(D_{i-1}, D_{i})$ falls into case $(+),(-)$ and $(\pm)$  respectively. Let $\cZ[D_{\bu}]$ be the substack of $\cZ_{\ol{k}}$ classifying 
\[
(\{x'_{i}\}\in X'^{r}; \{\cF_{i}\}\in \Hk^{r}_{U(n)}; \{\cL(D_{i})\xr{t'_{i}}\cF_{i}\}_{0\le i\le r})
\]
such that every $t'_{i}$ is saturated. Let 
\begin{equation}
\pi[D_{\bu}]: \cZ[D_{\bu}]\to  (X'_{\ol{k}})^{I_{0}}
\end{equation}
be the map recording those $x'_{i}$ for $i\in I_{0}$. Note that for $i\in I_{+}\cup I_{-}\cup I_{\pm}$, $x'_{i}$ is determined by $D_{\bu}$.


%Since $D_{\bu}$ determines $x'_{i}$ for those $i\in I_{+}\cup I_{-}\cup I_{\pm}$, the image of the projection $\cZ[D_{\bu}]_{\ol{\F}_q} \to (X'_{\ol{\F}_q})^{r}$ lands in a coordinate subspace $X'_{\ol{\F}_q}[D_{\bu}]$ that projects isomorphically to $(X'_{\ol{\F}_q})^{I_{0}}$. We denote the composition by $\pi[D_{\bu}]$
%\begin{equation}
%\pi[D_{\bu}]: \cZ[D_{\bu}]\to X'_{\ol{\F}_q}[D_{\bu}]\isom (X'_{\ol{\F}_q})^{I_{0}}.
%\end{equation}


\begin{prop}
\source{higher-siegel-weil-unitary-nonsingular-terms}
\notready\label{p:ShM dim} Let $n \geq 2$. \hfill
\source{higher-siegel-weil-unitary-nonsingular-terms}
\begin{enumerate}	
\item The substacks $\cZ[D_{\bu}]$ for $D_{\bu}\in \frD$ give a partition of $\cZ$.
\item Each geometric fiber of $\pi[D_{\bu}]$ has dimension $\le (n-1)|I_{+}|+(n-2)|I_{0}|$.
\item We have $\dim \cZ[D_{\bu}]\le r(n-1)$. The equality can only be achieved when  $I_{0}=\{1,2,\cdots, r\}$, i.e., all $D_{i}$ are equal to the same divisor of $X'$, which is then necessarily defined over $k$. \footnote{In this case, $\cZ[D_{\bu}]$ can be identified with the open substack $\mathring{\cZ}^{r}_{\cL(D_{0})}(a')\subset \cZ^{r}_{\cL(D_{0})}(a')$ (where $a'$ is the map $\cL(D_{0}) \rightarrow \sigma^* \cL(D_{0})^{\vee}$ induced from $a$) defined by requiring all the maps $t'_{i}: \cL(D_{0})\to \cF_{i}$ be saturated.}
\end{enumerate}
\end{prop}
\begin{proof}
(1) Each geometric point of $z\in \cZ$ defines a (unique) point $D_{\bu}\in \frD$ by taking the zero divisor of $t_{i}$, and then $z\in \cZ[D_{\bu}]$ by definition. 

(2) Let $\cH[D_{\bu}]$ be the substack of the fiber of $(\Hk^{r}_{\cM})_{\ol{k}}$ over $(\cL,a)\in \cA_{d}(k)$ classifying data $(\{x'_{i}\},\cL\xr{t_{i}}\cF_{i})$ such that $t_{i}$ extends to a map $t'_{i}: \cL(D_{i})\to \cF_{i}$ which is saturated. Note that  for $i\in I_{+}\cup I_{-}\cup I_{\pm}$, $x'_{i}$ is determined by $D_{\bu}$. Let $\cM[D_{i}]$ be the substack of the fiber of $\cM(1,n)_{\ol{k}}$ over $(\cL,a)\in \cA(1,n)(k)$ classifying maps $t:\cL\to \cF$ that extend to a saturated map $t': \cL(D_{i})\to \cF$. Then we have a Cartesian diagram of stacks over $\ol{k}$
\begin{equation}\label{ZHM}
\source{higher-siegel-weil-unitary-nonsingular-terms}
\xymatrix{\cZ[D_{\bu}]\ar[d]\ar[r] & \cH[D_{\bu}]\ar[d]^{(p_{0}, p_{r})}\\
\cM[D_{0}]\ar[r]^-{(\Id,\Frob)} & \cM[D_{0}]\times \cM[D_{r}].}
\end{equation}
Note since $D_{r}={}^{\t}D_{0}$, the Frobenius morphism sends $\cM[D_{0}]$ to $\cM[D_{r}]$. 

Let 
\begin{equation}
\Pi[D_{\bu}]: \cH[D_{\bu}]\to \cM[D_{r}]\times X'^{I_{0}}_{\ov k}
\end{equation}
be the projection $p_{r}$ and the map recording $x'_{i}$ for $i\in I^{0}$. 

\begin{claim}
\source{higher-siegel-weil-unitary-nonsingular-terms}
\notready The map $\Pi[D_{\bu}]$ is smooth and representable of relative dimension $(n-1)|I_{+}|+(n-2)|I_{0}|$.
\end{claim}

Assuming the claim, we finish the proof of (2). Indeed, we may apply Lemma \ref{l:VL} below to the Cartesian diagram \eqref{ZHM} to conclude, except that $\cM[D_{0}]$ is generally not a scheme. To remedy, we may restrict to a finite type open substack  $\cM[D_{0}]^{\le P}\subset \cM[D_{0}]$ by bounding Harder-Narasimhan polygon of $(\cF,h)$, and impose level structures on the Hermitian bundle $(\cF,h)$ at a closed point $x\in |X|$ to arrive at a scheme $\wt\cM[D_{0}]^{\le P}$ which is a torsor over $\cM[D_{0}]^{\le P}$ by an algebraic group $H$. Truncating and imposing the same level structures to $\cH[D_{\bu}]$ gives a  scheme $\wt\cH[D_{\bu}]^{\le P}$ (with legs away from $x$) such that $\wt\cH[D_{\bu}]^{\le P}/H=\cH[D_{\bu}]^{\le P}$. Let $\wt\cZ[D_{\bu}]^{\le P}$ be defined by a Cartesian diagram similar to \eqref{ZHM}, with $\cH[D_{\bu}]$ replaced by $\wt\cH[D_{\bu}]^{\le P}$ and $\cM[D_{i}]$ replaced by $\wt\cM[D_{i}]^{\le P}$ for $i=0,r$. We apply  Lemma \ref{l:VL} to conclude that the fibers of $\wt \cZ[D_{\bu}]^{\le P}\to X'^{I_{0}}_{\ol{k}}$ have dimension $\le (n-1)|I_{+}|+(n-2)|I_{0}|$. Now $\wt \cZ[D_{\bu}]^{\le P}/H(k)\isom \cZ[D_{\bu}]^{\le P}$ and for varying $P$ and $x$, $\cZ[D_{\bu}]^{\le P}$ cover $\cZ[D_{\bu}]$, hence the same dimension estimate holds for $\cZ[D_{\bu}]$.
%
%By Lemma \ref{l:VL}, to show that $\pi[D_{\bu}]$ is smooth of relative dimension $(n-1)|I_{+}|+(n-2)|I_{0}|$, it suffices to show that 
%\begin{equation}
%\Pi[D,I]: \cH[D_{\bu}]\to \cM[D_{r}]\times X'^{I_{0}}
%\end{equation}
%is smooth of the same relative dimension. 

It remains to prove the claim.

%$\{x'_{i}\}_{i\in I_{0}}\in {X'_{\ol{k}}}^{I_{0}}$ 
For $r\ge j\ge 0$, let $\cH_{\ge j}$ be the moduli stack defined similarly to $\cH[D_{\bu}]$ but classifying saturated maps $\{t_{i}: \cL(D_{i}) \to \cF_{i}\}_{j\le i\le r}$ (lying over $a$) only for $i$ in the indicated range. We can factorize $\Pi[D_{\bu}]$ as
\begin{equation}
\Pi[D_{\bu}]: \cH[D_{\bu}]=\cH_{\ge0} \xr{ \Pi_{1}} \cH_{\ge1}\xr{\Pi_{2} }\cdots \xr{ \Pi_{r}}\cH_{\ge r}=\cM[D_{r}]\times X'^{I_{0}}_{\ov k}.
\end{equation}
The desired smoothness and relative dimension claims would follow from the following four statements:
\begin{enumerate}
\item[($H0$)] If $i\in I_{0}$, then $\Pi_{i}$ exhibits $\cH_{\ge i-1}$ as an open substack in a $\PP^{n-2}$-bundle over $\cH_{\ge i}$.
\item[($H+$)] If $i\in I_{+}$, then $\Pi_{i}$ exhibits $\cH_{\ge i-1}$ as an open substack in a $\PP^{n-1}$-bundle over $\cH_{\ge i}$.
\item[($H-$)] If $i\in I_{-}$, then $\Pi_{i}$ is an isomorphism.
\item[($H\pm$)] If $i\in I_{\pm}$, then $\Pi_{i}$ is an open immersion.
\end{enumerate}
We next establish each of these statements. 

\emph{Proof of }($H0$). When $i\in I_{0}$, $D_{i-1}=D_{i}$. We write the modification $\cF_{i-1}\dashrightarrow\cF_{i}$ as 
\begin{equation}\label{lower middle}
\source{higher-siegel-weil-unitary-nonsingular-terms}
\begin{tikzcd}
\cF_{i-1} & \cF^{\flat}_{i-1/2}\ar[r, "\s(x'_{i})", hook] \ar[l, "{x'_{i}}"', hook']  & \cF_{i} 
\end{tikzcd}
\end{equation}
Here both arrows have cokernel of length one supported at the labelled points. Such modifications of $\cF_{i}$ are parametrized by a hyperplane $H$ in the fiber $\cF_{i}|_{\s(x'_{i})}$. 
The requirement that $t_{i}:\cL(D_{i})\to \cF_{i}$ should land in $\cF^{\flat}_{i-1/2}$ is equivalent to the (closed) condition that $H$ should contain the line given by the image of $\cL(D_{i})|_{\s(x'_{i})}$. This cuts out a $\PP^{n-2}$ in the space of hyperplanes $H \subset \cF_{i}|_{\s(x'_{i})}$. The further requirement that $t_{i-1}:\cL(D_{i})\to \cF^{\flat}_{i-1/2}\to \cF_{i-1}$ be saturated is an open condition. 

This argument globalizes in the evident way, exhibiting that $\Pi_i$ as an open substack in a $\PP^{n-2}$-bundle. This applies similarly for the analogous arguments below for the other cases. 


\emph{Proof of} ($H+$). When $i\in I_{+}$, we have $D_{i-1}=D_{i}-\s(x'_{i})$. We write the modification of $\cF_{i}$ as in \eqref{lower middle}. This time the choice of the $\cF_{i}\leftarrow\cF^{\flat}_{i-1/2}$ is the open subset of those hyperplanes $H\in \PP(\cF_{i}|_{\s(x'_{i})})$ that do not contain the image of $\cL(D_{i})|_{\s(x'_{i})}$. The requirement that $\cL(D_{i-1})=\cL(D_{i}-\s(x'_{i}))\to \cF^{\flat}_{i-1/2}\to \cF_{i-1}$ be saturated at $x'_{i}$ imposes a further open condition. 

\emph{Proof of} ($H-$). When $i\in I_{-}$, we have $D_{i-1}=D_{i}+x'_{i}$. We write the modification as 
\begin{equation}\label{upper middle}
\source{higher-siegel-weil-unitary-nonsingular-terms}
\begin{tikzcd}
\cF_{i-1}\ar[r, "\s(x'_{i})", hook] & \cF^{\sh}_{i-1/2} & \cF_{i}\ar[l, "x'_{i}"', hook' ]
\end{tikzcd} 
\end{equation}
where both arrows have cokernel of length one supported at the labelled points. Now $t_{i}: \cL(D_{i})\to \cF_{i}$ is required to extend to $\cL(D_{i}+x'_{i})\to \cF^{\sh}_{i-1/2}$. This determines the upper modification $\cF_{i}\to \cF^{\sh}_{i-1/2}$ uniquely, which in turn determines the lower modification $\cF^{\sh}_{i-1/2}\leftarrow \cF_{i-1}$ as well. We get a map $t'_{i-1}: \cL(D_{i-1})=\cL(D_{i}+x'_{i})\to \cF^{\sh}_{i-1/2}$. We claim that $t'_{i-1}$ automatically lands in $\cF_{i-1}$. Indeed, the claim is equivalent to saying that under the pairing between $\cF_{i}|_{x'_{i}}$ and $\cF_{i}|_{\s(x'_{i})}$, the images of $t_{i}(x'_{i})$ and $t_{i}(\s(x'_{i}))$ pair to zero. The latter statement is equivalent to saying that the induced Hermitian map $a'_{i}: \cL(D_{i})\to \s^{*}(\cL(D_{i}))^{\vee}$ vanishes at $x'_{i}$. But we know that the divisor of $a'_{i}$ is $\nu^{*}D_{a}-D_{i}-\s(D_{i})$. Since $D_{i-1}=D_{i}+x'_{i}$ satisfies $\nu(D_{i-1})\le D_{a}$ by assumption, we see that $\nu^{*}D_{a}-D_{i}-\s(D_{i})\ge x'_{i}+\s(x'_{i})$, hence $a'_{i}$ is guaranteed to vanish at $x'_{i}$ and $\s(x'_{i})$. This shows that there is a unique lifting of any point of $\cH_{\ge i}$ to $\cH_{\ge i-1}$, hence $\Pi_{i}$ is an isomorphism.  

\emph{Proof of} ($H\pm$). When $i\in I_{\pm}$, we have $D_{i}+x'_{i}=D_{i-1}+\s(x'_{i})$. We write the modification as in \eqref{upper middle}. As in the case ($H-$), the requirement that $t_{i}: \cL(D_{i})\to \cF_{i}$ should extend to $\cL(D_{i}+x'_{i})\to \cF^{\sh}_{i-1/2}$ determines the modification. Then we automatically get a map $t_{i-1}: \cL(D_{i-1})=\cL(D_{i}+x'_{i}-\s(x'_{i}))\to \cF_{i-1}$; the requirement that $t_{i-1}$ be saturated is an open condition. Therefore $\Pi_{i}$ is an open immersion in this case.

(3) By (2) we have 
\[
\dim \cZ[D_{\bu}]\le (n-1)|I_{+}|+(n-2)|I_{0}|+|I_{0}|=(n-1)(|I_{+}\cup I_{0}|)\le r(n-1).
\]
Equality holds only if $I_{-}$ and $I_{\pm}$ are empty. However, by degree reasons we have $|I_{+}|=|I_{-}|$, so in the equality case we must have $I_{+}=\vn$ as well. We conclude that equality can only be achieved if $I_{0}=\{1,2,\cdots, r\}$; in other words, all $D_{i}$ must be the same. In particular, since $D_r = \ft D_0$, this forces $D_0$ to be defined over $k$. 
\end{proof}

The Lemma below, a slight variant of \cite[Lemma 2.13]{Laff18}, was used above.

\begin{lemma}[Variant of {\cite[Lemma 2.13]{Laff18}}]
\source{higher-siegel-weil-unitary-nonsingular-terms}
\notready\label{l:VL}
\source{higher-siegel-weil-unitary-nonsingular-terms}
Let $W,Z, T$ be schemes of finite type over $\ov{k}$. Let $Z^{(1)}$ be the Frobenius twist of $Z$ (i.e., the pullback of $Z$ under the $q$-Frobenius $\Spec \ov{k}\to \Spec \ov{k}$). Let $\wt h = (h_1, h_T) \co W \rightarrow Z^{(1)} \times T$ be smooth of relative dimension $d$, and $h_{0}\co W \rightarrow Z$ be an arbitrary map. Define $V$ as the fibered product
\[
\begin{tikzcd}
V \ar[r] \ar[d] & W \ar[d, "{(h_{0},h_{1})}"] \\
Z \ar[r, "{(\Id, \Frob)}"] & Z \times Z^{(1)}
\end{tikzcd}
\]
Then  each fiber of the composition map $V \rightarrow W\xr{h_{T}}T$ has dimension $\le d$. 
\end{lemma}

\begin{proof} Restricting $W$ over a point $t\in T(\ov{k})$, we reduce to the case $T$ itself is the point $\Spec\ov{k}$. We may assume $Z=\Spec R$ where $R=\ov k[x_{1},\ldots, x_{l}]/I$. Let $R^{(1)}=\ov k\ot_{\ov k}R\cong \ov k[\xi_{1},\ldots, \xi_{l}]/I^{(1)}$ be the base change of $R$ under $\Frob_{q}$, where $\xi_{i}=1\ot x_{i}$ . Since $h_{1}: W\to Z^{(1)}$ is smooth of relative dimension $d$,  by Zariski localizing we may assume $W=\ov k[\xi_{1},\ldots, \xi_{l}, y_{1},\ldots, y_{m+d}]/(I^{(1)}, r_{1},\ldots, r_{m})$,  with $(\pderiv{r_i}{y_j})_{j=1}^m$ having rank $m$ ($r_{i}\in \ov k[\xi_{1},\ldots, \xi_{l}, y_{1},\ldots, y_{m+d}]$). Under $h_{0}: W\to Z$, the coordinates $x_{i}$ of $Z$ pullback to functions $\ov f_{i}$ on $W$, $1\le i\le l$. We lift $f_{i}$ to polynomials $f_{i}\in \ov k[\xi_{1},\ldots, \xi_{l}, y_{1},\ldots, y_{m+d}]$. 

By definition, $V$ has the form 
\begin{eqnarray*}
%V &\cong& \Spec \dfrac{ \ov k[x_{1},\cdots, x_{l}, \xi_{1},\ldots, \xi_{l}, y_1, \ldots, y_m, y_{m+1} \ldots, y_{m+d}  ]}{(I(x), I^{(1)}(\xi), x_{1}-f_1, \ldots, x_{l}-f_{l}, \xi_{1}-x^{q}_{1},\ldots, \xi_{l}-x_{l}^{q}, r_1, \ldots, r_m)}\\
V & \cong&  \Spec\dfrac{ \ov k[\xi_{1},\ldots, \xi_{l}, y_1, \ldots, y_m, y_{m+1} \ldots, y_{m+d}  ]}
{(I^{(1)}(\xi), g_{1},\ldots, g_{l}, r_1, \ldots, r_m)}
\end{eqnarray*}
where $g_{i}=\xi_{i}-f_{i}^{q}$. In particular,  $V$ is a closed subscheme of 
\begin{equation}
U:=\Spec\Big( \ov k[\xi_{1},\ldots, \xi_{l}, y_1, \ldots, y_m, y_{m+1} \ldots, y_{m+d}  ]/ (g_{1},\ldots, g_{l}, r_1, \ldots, r_m)\Big).
\end{equation}
The Jacobian matrix for the defining equations of $U$ has the form 
\[
\begin{pmatrix}
\boxed{\pderiv{g_i}{\xi_j}}_{j=1}^l  & \boxed{\pderiv{g_i}{y_j}}_{j=1}^m & \boxed{\pderiv{g_i}{y_j}}_{j=m+1}^{m+d}  \\
\boxed{\pderiv{r_i}{\xi_j}}_{j=1}^l  & \boxed{\pderiv{r_i}{y_j}}_{j=1}^m & \boxed{\pderiv{r_i}{y_j}}_{j=m+1}^{m+d} 
\end{pmatrix} = \begin{pmatrix} \Id_l & 0 &  0 \\ * & \text{invertible}_m & *  \end{pmatrix},
\]
which evidently has rank $l+m$. Hence $U$ is smooth of dimension $d$. Since $V\incl U$, $\dim V\le d$.
\end{proof}



%
%\begin{lemma}[Variant of {\cite[Lemma 2.13]{Laff18}}]\label{l:VL}
%Let $W,Z,T$ be schemes of finite type over $\ov{k}$ with $Z$ smooth. Let $Z^{(1)}$ be the Frobenius twist of $Z$ (i.e., the pullback of $Z$ under the $q$-Frobenius $\Spec \ov{k}\to \Spec \ov{k}$). Let $\wt h = (h_1, h_T) \co W \rightarrow Z^{(1)} \times T$ be smooth of relative dimension $d$, and $h_{0}\co W \rightarrow Z$ be an arbitrary map. Define $V$ as the fibered product
%\[
%\begin{tikzcd}
%V \ar[r] \ar[d] & W \ar[d, "{(h_{0},h_{1})}"] \\
%Z \ar[r, "{(\Id, \Frob)}"] & Z \times Z^{(1)}
%\end{tikzcd}
%\]
%Then the map $V \rightarrow W\xr{h_{T}}T$ is smooth of relative dimension $d$. 
%
%In particular, if $W,Z,T$ are equidimensional, then $V$ is equidimensional and its dimension agrees with its ``expected dimension'' $\dim W - \dim Z$. 
%\end{lemma}
%
%\begin{proof}
%We make the same initial reductions as in \cite[Proof of Lemma 2.13]{Laff18}. In particular, by \'{e}tale localization on $Z$, we may assume that $Z = \A^l$. Then, by Zariski localization on $T$ and $W$, we may assume that $T$ is affine with coordinate ring $\Cal{O}(T)$, and $W$ is open in $\Spec \left( \Cal{O}(T)[x_1, \ldots, x_l, y_1, \ldots, y_{m+d}]/(r_1, \ldots, r_m) \right)$ with $(\pderiv{r_i}{y_j})$ having rank $m$. 
%
%Then $U$ has the form 
%\[
%U \cong \Spec \Big( \Cal{O}(T)[x_1, \ldots, x_l, y_1, \ldots, y_m, y_{m+1} \ldots, y_{m+d}  ] / (g_1, \ldots, g_l, r_1, \ldots, r_m)\Big),
%\]
%where $g_i = x_i - h_1(x_i)^q$. Therefore, the Jacobian of the ideal $(g_1, \ldots, g_l, r_1, \ldots, r_m)$ has the form 
%\[
%\begin{pmatrix}
%\boxed{\pderiv{g_i}{x_j}}_{j=1}^l  & \boxed{\pderiv{g_i}{y_j}}_{j=1}^m & \boxed{\pderiv{g_i}{y_j}}_{j=m+1}^{m+d}  \\
%\boxed{\pderiv{r_i}{x_j}}_{j=1}^l  & \boxed{\pderiv{r_i}{y_j}}_{j=1}^m & \boxed{\pderiv{r_i}{y_j}}_{j=m+1}^{m+d} 
%\end{pmatrix} = \begin{pmatrix} \Id_l & 0 &  0 \\ * & \text{invertible}_m & *  \end{pmatrix},
%\]
%which evidently has rank $l+m$. Hence $U \rightarrow T$ is smooth of relative dimension $d$. 
%\end{proof}



%%%%%%%%COMMENTING OUT%%%%%%%5
\begin{comment}
\mnote{Y: since we deleted the smoothness statement for $\cM(m,n)|_{\cA(m,n)}$, the following proof is not self-contained. Either delete this Cor. or add back the smoothness of $\cM(m,n)|_{\cA(m,n)}$.}

\begin{cor}
\source{higher-siegel-weil-unitary-nonsingular-terms}
\notready\label{cor: KR is LCI}
\source{higher-siegel-weil-unitary-nonsingular-terms}
\begin{enumerate}
\item For any $n \geq 1$, a line bundle $\cL$ on $X'$ and $a\in \cA_{\cL}(k)$, the stack $\cZ_{\cL}^{r}(a)$ is LCI of pure dimension $r(n-1)$. 
\item When $n\ge2$, for any effective divisor $D$ on $X'$ such that $\nu(D)\le D_{a}$, the locally closed substack $\mathring{\cZ}_{\cL(D)}^{r}(a)\subset \cZ_{\cL}^{r}(a)$ where $t_{i}$ extends to a {\em saturated} map $\cL_{i}(D)\to \cF$ for all $0\le i\le r$ has pure dimension $r(n-1)$, and $\cup_{\nu(D)\le D_{a}}\mathring{\cZ}_{\cL(D)}^{r}(a)$ is dense in $\cZ^{r}_{\cL}(a)$.
\end{enumerate}
%\[
%\Sht^{r}_{\cM}=\coprod_{\cL} \coprod_{a \in \cA_{\cL}^{\all}(k)} \cZ_{\cL}^{r}(a)^*
%\]
%has pure dimension $r(n-1)$.
\end{cor}
\begin{proof} 
(1) In this proof we write $\cM=\cM(1,n)|_{\cA(1,n)}$. By Proposition \ref{} $\cM$ is smooth. We can rewrite $\Sht^{r}_{\cM}$ by the Cartesian diagram
\begin{equation}\label{eq: shtuka cycle class} 
\source{higher-siegel-weil-unitary-nonsingular-terms}
\begin{tikzcd}
 \Sht^{r}_{\cM}\ar[r]\ar[d]      &  (\Hk^{1}_{\cM})^{r}\ar[d, "{(p_{0},p_{1})^r}"] \\
\cM^{r} \ar[r, "\varphi"] &  \cM^{2r}
\end{tikzcd}
\end{equation}
where $\ph(\xi_{0},\cdots, \xi_{r-1})=(\xi_{0},\xi_{1},\xi_{1},\xi_{2},\xi_{2},\cdots, \xi_{r-1}, \Frob(\xi_{0}))$, and $(p_{0},p_{1}): \Hk^{1}_{\cM}\to \cM\times \cM$ are the usual projections. 

Note that $\Sht^{r}_{\cM}=\coprod_{(\cL,a)\in \cA(1,n)(k)}\cZ^{r}_{\cL}(a)$. By Lemma \ref{lem: dim Hk_M},  $\Sht^{r}_{\cM}$ has local dimension $\ge r(n-1)$ everywhere, therefore the same is true for each $\cZ^{r}_{\cL}(a)$. If $n \geq 2$, then by Proposition \ref{p:ShM dim}(3) the equality holds. 

If $n = 1$, then by Proposition \ref{prop: KR cycle proper}, $\cZ_{\cL}^r(a)$ is supported on the the fiber of $\Sht^r_{U(1)}$ over finitely many points in $X'^{r}$ (the legs must appear in $D_{a}$), which is 0-dimensional by Lemma \ref{lem: shtuka geometry}. So again the equality of dimensions holds. 

%then we have 
%\[
%\cZ_{\cL}^{r} = \bigcup_{a \in \cA_{\cL}(k)} \cZ_{\cL}^r(a)
%\]
%and 

(2) Since $\mathring{\cZ}_{\cL(D)}^{r}(a)$ is open in $\cZ_{\cL(D}^{r}(a)$, which is of pure dimension $r(n-1)$ by part (1), the same is true for $\mathring{\cZ}_{\cL(D)}^{r}(a)$. By Proposition \ref{p:ShM dim}(3),  the complement of $\cup_{\nu(D)\le D_{a}}\mathring{\cZ}_{\cL(D)}^{r}(a)$ is the union of the  strata $\cZ[D_{\bu}]$ where $D_{\bu}$ is not constant, and the latter has dimension strictly less than $r(n-1)$. Since $\cZ^{r}_{\cL}(a)$ is LCI, the union of $\mathring{\cZ}_{\cL(D)}^{r}(a)$ must be dense.
\end{proof}
\end{comment}
%%%%%%%%END OF COMMENTED PART


\subsection{The case $m=1$ and $a=0$}\label{ssec: m=1 a=0 case}
\source{higher-siegel-weil-unitary-nonsingular-terms}
We keep the notations from \S\ref{ss:strat KR div}. In this subsection we extend the discussion in \S\ref{ss:strat KR div} to the case $a=0$. Fix a line bundle $\cL\in \Pic(X')$. Recall from Definition  \ref{def:Z(0)} that $\cZ_{\cL}^{r}(0)^{\circ}$ is the moduli stack classifying hermitian shtukas $(\{x'_{i}\},\{\cF_{i}\})$ together with compatible maps $\{\cL\xr{t_{i}}\cF_{i}\})$ with $t_{i}$ injective (fiberwise over the test scheme $S$) and the image of $t_{i}$ being isotropic.  In this subsection, let
\begin{equation}
\cZ:=\cZ_{\cL}^{r}(0)^{\circ}_{\ov k}.
\end{equation}
If $n=1$ then $\cZ_{\cL}^{r}(0)^{\circ}=\vn$. We always assume $n\ge2$ below. 

\subsubsection{Indexing set for strata} Let $I_{0}\sqcup I_{+}\sqcup I_{-}$ be a partition of $\{1,2,\cdots, r\}$ such that $|I_{+}|=|I_{-}|$.  We denote this partition simply by $I_{\bu}$. For any $N\in \Z_{\ge0}$, define $\frD(N;I_{\bu})$ to be the moduli space of sequences of effective divisors $(D_{i})_{0\le i\le r}$ on $X'_{\ov k}$ such that
\begin{enumerate}
\item $\deg(D_{0})\le N$.
\item For $?=0,+$ or $-$, and $i\in I_{?}$, the pair $(D_{i-1},D_{i})$ belongs to the corresponding Case (?) listed in the beginning of \S\ref{ss:strat KR div} (for some $x'_{i}\in X'_{\ov k}$ in the case $?=+$ or $-$).
\item $D_{r}={}^{\t}D_{0}$. 
\end{enumerate}
We have a map recording the points $x'_{i}$ for $i\in I_{+}\cup I_{-}$:
\begin{equation}
(\pi_{+}, \pi_{-}): \frD(N;I_{\bu})\to (X'_{\ol{k}})^{I_{+}}\times (X'_{\ol{k}})^{I_{-}}.
\end{equation}

\begin{lemma}
\source{higher-siegel-weil-unitary-nonsingular-terms}
\notready\label{l:dim D}
\source{higher-siegel-weil-unitary-nonsingular-terms}
The map $\pi_{+}: \frD(N;I_{\bu})\to (X'_{\ol{k}})^{I_{+}}$ is quasi-finite.
\end{lemma}
\begin{proof}
For a fixed geometric point $(x'_{i})_{i\in I_{+}} \in (X'_{\ol{k}})^{I_{+}}(\ol{k})$, its fiber in $\frD(N;I_{\bu})$ consists of $(D_{0}, \{x'_{i}\}_{i\in I_{-}})$ such that $\deg D_{0}\le N$ and 
\begin{equation}\label{eqn D0}
\source{higher-siegel-weil-unitary-nonsingular-terms}
D_{0}+\sum_{i\in I_{+}} \s(x'_{i})={}^{\t}D_{0}+\sum_{i\in I_{-}} x'_{i}.
\end{equation}
Let $v\in |X'|$ be a closed point that intersects $\supp D_{0}$. If $\deg(v)>N$, then $D_{0}$ cannot contain all geometric points over $v$ and hence there exists a geometric point $y|v$ such that $y\in \supp {}^{\t}D_{0}$ but $y\notin \supp D_{0}$. By \eqref{eqn D0}, $y=\s(x'_{i})$ for some $i\in I_{+}$.  Therefore points in $D_{0}$ are either over closed points of degree $\le N$, or in the Galois orbit of $\s(x'_{i})$ for some $i\in I_{+}$. This leaves finitely many possibilities for $D_{0}$, hence for $\{x'_{i}\}_{i\in I_{-}}$ as well.
\end{proof}



%Let $U=X'-\cup_{m\le N}X'(\F_{q^{m}})$. Then the above equation implies that $\sum_{i\in I_{+}} \s(x'_{i})$ and $\sum_{i\in I_{-}} x'_{i}$ are equal when restricted to $U$ (the only part that can interact with $D_{0}$ and ${}^{\t}{D_{0}}$ are points of degree $\le N$, for $\deg(D_{0})=N$). Therefore, the points $\{x'_{i}\}_{i\in I_{-}}$ either have degree $\le N$, or already appear in $\{\s(x'_{i})\}_{i\in I_{+}}$. This leaves finitely many choices for $\{x'_{i}\}_{i\in I_{-}}$. Once $\{x'_{i}\}_{i\in I_{-}}$  is fixed, there are also only finitely many possible choices for $D_{0}$ satisfying the equation \eqref{eqn D0}. \mnote{TF: isn't this finiteness true even without fixing $\{x'_{i}\}_{i\in I_{-}}$, since $D_0$ was effective of some bounded degree?} Therefore the fiber of $\frD(N;I_{\bu})$ over $(x'_{i})_{i\in I_{+}}$ is finite.


\subsubsection{Definition of strata} For a partition $I=(I_{0},I_{+},I_{-})$ of $\{1,2,\cdots, r\}$, define $\cZ[N;I_{\bu}]$ to be the stack classifying 
\[
(\{D_{i}\}_{0\le i\le r},( \{x'_{i}\}_{1\le i\le r}, \{\cF_{i}\}_{0\le i\le r})\in \Hk^{r}_{U(n)}, \{\cL\xr{t_{i}}\cF_{i}\}_{0\le i\le r})
\]
such that $\{D_{i}\}\in \frD(N;I_{\bu})$ with image $\{x'_{i}\}_{i\in I_{?}}$ under $\pi_{?}$ ($?=+,-$), and $t_{i}$ extends to a saturated embedding $\cL(D_{i})\inj \cF_{i}$. We have a map
\begin{equation}
\pi[N;I_{\bu}]: \cZ[N;I_{\bu}]\to (X'_{\ol{k}})^{I_{0}}\times \frD(N;I_{\bu}).
\end{equation}
The following is the analog of Proposition \ref{p:ShM dim} when $a=0$.

\begin{prop}
\source{higher-siegel-weil-unitary-nonsingular-terms}
\notready\label{p:dim ZL0} Let $n \geq 2$. \hfill 
\source{higher-siegel-weil-unitary-nonsingular-terms}
\begin{enumerate}
\item For varying $N\in \Z_{\ge0}$ and partitions $I_{\bu}$ of $\{1,2,\cdots, r\}$ such that $|I_{+}|=|I_{-}|$, the substacks  $\cZ[N;I_{\bu}]$ give a partition of $\cZ$. 
\item The fibers of the map $\pi[N;I_{\bu}]$ have dimension $\le(n-1)|I_{+}|+(n-2)|I_{0}|$.
\item We have $\dim \cZ[N;I_{\bu}]\le r(n-1)$. Moreover, when $n\ge3$, the equality  can only be achieved when $I_{0}=\{1,2,\cdots, r\}$, i.e., all $D_{i}$ are equal to the same divisor of $X'$ defined over $k$.
\end{enumerate}
\end{prop}
\begin{proof}
(1) is similar to Proposition \ref{p:ShM dim}(1), except we have to argue that the strata are only non-empty for $|I_+|=|I_-|$, and that Case $(\pm)$ cannot appear for points in $\cZ=\cZ^{r}_{\cL}(0)^{\circ}$. The first statement follows from the assumption that $D_r = \ft D_0$ has the same degree as $D_0$. For the second statement, suppose Case $(\pm)$ happens for the modification 
\[
\cF_{i-1}\leftarrow \cF_{i-1/2}^{\flat}\to \cF_{i},
\]
and let $H\subset \cF_{i-1,x'_{i}}$ be the hyperplane that is the image of $\cF_{i-1/2}^{\flat}$. Then $H^{\bot}\subset \cF_{i-1,\s(x'_{i})}$ is the line along which the upper modification $\cF_{i-1/2}^{\flat}\inj \cF_{i}$ is performed. Let $\ell_{x'_{i}}$ (resp. $\ell_{\s(x'_{i})}$) be the image of $\cL(D_{i-1})\to \cF_{i-1}$ at $x'_{i}$ (resp. at $\s(x'_{i})$). Since the image of $\cL(D_{i-1})$ is isotropic (because $a=0$), $(\ell_{x'_{i}},\ell_{\s(x'_{i})})=0$ under the pairing between $\cF_{i-1,x'_{i}}$ and $\cF_{i-1,\s(x'_{i})}$. The condition $D_{i}+x'_{i}=D_{i-1}+\s(x'_{i})$ happens only if $\ell_{x'_{i}}\not\subset H$ and $\ell_{\s(x'_{i})}=H^{\bot}$. This contradicts the fact that $(\ell_{x'_{i}},\ell_{\s(x'_{i})})=0$.

(2) is proved in the same way as Proposition \ref{p:ShM dim}(2). 

(3) Applying (2) and Lemma \ref{l:dim D} we get
\begin{eqnarray*}
\dim\cZ[N;I_{\bu}]&=&  (n-1)|I_{+}|+(n-2)|I_{0}|+|I_{0}|+\dim \frD(N;I_{\bu})\\
&\le & (n-1)|I_{+}|+(n-1)|I_{0}|+|I_{+}|=(n-1)|I_{0}|+n|I_{+}|.
\end{eqnarray*}
Since $n\ge2$,  we have $n\le2(n-1)$, therefore the above is $\le (n-1)(|I_{0}|+2|I_{+}|)=r(n-1)$. When $n\ge3$, we have strict inequalities $n<2(n-1)$, so equality can only be achieved when $I_{+}$, hence $I_{-}$, are all empty.
\end{proof}




%\subsection{Dimension estimates}
